\section{はじめに}
近年，強化学習を実運用へ適用するためには，環境が変化しても短時間で性能を回復・向上できる適応性が重要である．
強化学習は，エージェントが試行錯誤を通じて最適な行動方策を獲得する枠組みであり
\cite{SuttonBarto1998,WatkinsDayan1992}，
ゲームやロボティクスに加え，物流・製造計画など多様な領域での活用が期待されている．
一方で，反復的な探索と学習には多大な試行回数と計算時間を要し，
実環境では学習効率の低さが主要な課題となる
\cite{Thrun1995}．
例えば倉庫の入出庫スケジュール最適化や工場のジョブショップスケジューリングでは，
設備配置や需要が変化するたびにゼロから学習し直すことは現実的でない場合が多い．

この課題に対する有力なアプローチとして転移学習がある
\cite{TaylorWhitesonStone2007}．
転移学習では，変化前の環境（ソース）で獲得した知識を新しい環境（ターゲット）の学習に再利用することで，
学習初期の性能向上や総学習時間の短縮が期待できる．
とりわけ，方策や価値関数の再利用は環境構造が類似する場合に効果的であるが，
ソース知識をどの程度信頼して利用するかは容易ではなく，
不適切な再利用はターゲット適応を阻害しうる
\cite{TaylorKuhlmannStone2008}．

過去に学習した方策を探索へ組み込む枠組みとして，Fern\'andez と Veloso は
方策ライブラリからエピソード単位で方策を選択し，確率的に再利用する手法を提案している
\cite{Fernandez2006}．
また，学習を改善するために追加知識を探索へ注入する研究として，
助言ルールによる学習支援
\cite{MaclinShavlik1996}，
模倣に基づく方策獲得
\cite{NicolescuMataric2003}，
階層化による知識再利用
\cite{BartoMahadevan2003,KonidarisBarto2007}
などが報告されている．

本研究では，転移学習の枠組みにおける方策選択手法の影響を体系的に評価することを目的とする．
具体的には，既存研究で用いられているボルツマン選択に加え，
$\varepsilon$-greedy および UCB を導入し，
3種類のバンディットアルゴリズムに基づく方策選択を同一条件下で比較実験する．
さらに，報酬設計の違いが比較結果に与える影響にも着目し，
従来の報酬に加えて，各ステップでのマイナス報酬および壁衝突時のペナルティを導入した報酬設計においても同様の検証を行う．
これにより，探索と活用のバランスが異なる方策選択が，
報酬構造の変化に対してどのように挙動を変えるのかを明らかにし，
実運用を見据えた転移強化学習の設計指針を得ることを目指す．

